Parameter-efficient fine-tuning (PEFT) has become a central approach for adapting large pretrained language models while updating only a small fraction of their parameters. Low-Rank Adaptation (LoRA)~\citep{hu2022lora} is among the most widely used PEFT methods, showing that additive low-rank updates can achieve strong downstream performance with substantially lower memory and parameter cost than full fine-tuning. Beyond efficiency, LoRA also provides an important empirical insight: fine-tuning updates tend to amplify task-relevant singular directions that are already present in the pretrained weights. We refer to this view as the \emph{pretraining-finetuning paradigm}: pretraining establishes a rich set of representational directions, while fine-tuning primarily recalibrates their task-specific importance.

This perspective suggests that the singular structure of pretrained weights is not merely a tool for compression or initialization, but a meaningful coordinate system for adaptation. Recent PEFT methods have increasingly moved in this direction. AdaLoRA, VeRA, MiLoRA, and PiSSA~\citep{zhang2023adalora, kopiczko2023vera, wang2025milora, meng2024pissa} use low-rank, spectral, or SVD-inspired structure to improve rank allocation, initialization, or parameter efficiency. However, many existing methods use the pretrained spectrum indirectly, for example through initialization, random bases, or soft allocation rules. They do not explicitly analyze how the adapter update decomposes in the pretrained singular coordinates throughout training.

In this work, we study PEFT directly in the singular coordinate system of the pretrained weight matrices. We distinguish between two levels of the proposed construction. First, we introduce idealized spectral adapters that make the preservation-adaptation structure explicit. Second, we introduce their practical unit-invariant variants, which add lightweight diagonal coordinate scalers.

We begin with a structurally transparent adapter, \textbf{LinLoRA}, which freezes the pretrained singular vectors and learns only singular-value coefficients in the complementary tail block. In this idealized form, the pretrained and fine-tuned components occupy orthogonal blocks of the fixed pretrained basis, giving a direct mechanism for separating preserved information from injected task-specific information.

This strict structure also exposes a limitation. Because LinLoRA adapts within a fixed representational basis, its effectiveness depends on whether the prescribed tail basis aligns with the dominant directions of the optimal task-specific perturbation. We formalize this as a \emph{fixed-basis adaptation bottleneck}: under a PAC-Bayesian perturbation perspective, a fixed-basis adapter can reduce the residual error bound only to the extent that its adaptation subspace captures the relevant task directions. If this alignment fails, the adapter may leave the leading components of the optimal perturbation uncaptured.

To address this bottleneck while preserving parameter efficiency, we introduce \textbf{OrthoLoRA}. Rather than learning dense new basis matrices, OrthoLoRA freezes the pretrained tail singular vectors and learns compact orthogonal rotations \(R_U\) and \(R_V\) within the selected tail subspace. This allows the adapter to rotate its basis toward task-relevant directions while keeping the parameter cost tied to the spectral budget rather than to the full input and output dimensions. In particular, the additional trainable cost scales as \(\mathcal{O}(k_{vec}^2)\), rather than \(\mathcal{O}(d_{out}k_{vec}+d_{in}k_{vec})\).

We then introduce the unit-invariant variants of these adapters by adding trainable diagonal scalers \(E\) and \(D\), inspired by the unit-invariant SVD perspective of \citet{uhlmann2018generalized}. The prefix \textbf{UI} denotes this unit-invariant extension. Thus, \textbf{UILinLoRA} is the unit-invariant version of LinLoRA: it keeps the fixed-basis, magnitude-only spectral structure, but adds the diagonal scalers. Similarly, \textbf{UIOrthoLoRA} is the unit-invariant version of OrthoLoRA: it combines compact orthogonal rotations with the same diagonal scalers.

The diagonal scalers relax the strict confinement of pure LinLoRA and OrthoLoRA by allowing unit-level coordinate adaptation around the pretrained tail subspace. This added flexibility is important in practical settings, where task-specific perturbations may not be perfectly aligned with the original pretrained singular geometry. At the same time, the relaxation is not unstructured: we analyze the UIOrthoLoRA update in pretrained singular coordinates and derive explicit bounds on the scaler-induced mixing between preserved leading components and adapted tail components.

Empirically, the main paper evaluates UILinLoRA and UIOrthoLoRA on GLUE using RoBERTa-base and RoBERTa-large results. On RoBERTa-base, UIOrthoLoRA achieves the highest average score among the compared VeRA, LoRA, and RandLoRA baselines, reaching 85.5 with 0.4M trainable parameters. UILinLoRA also remains strong, reaching 84.68 with only 0.08M trainable parameters, highlighting the efficiency of spectral magnitude adaptation. Additional generation and retention-oriented experiments are reported in the appendix as supporting evidence that the proposed spectral parameterization remains expressive while preserving the paper's main preservation-adaptation motivation.
Our main contributions are:
\begin{itemize}
    \item \textbf{Spectrum-aware PEFT and the fixed-basis bottleneck.}
    We formulate adapter updates in the singular coordinate system of pretrained weights, introduce LinLoRA as a fixed-basis spectral adapter, and show that its effectiveness depends on alignment between the tail basis and the dominant task-specific perturbation directions.

    \item \textbf{Compact orthogonal adaptation with controlled mixing.}
    We introduce OrthoLoRA, which adds compact tail-subspace rotations, and UIOrthoLoRA, which further adds unit-invariant diagonal scalers. We derive bounds on the resulting leading-tail mixing induced by the scalers.

    \item \textbf{Empirical validation.}
    We show strong parameter efficiency on RoBERTa-base and RoBERTa-large on GLUE,
    report additional generation and retention-oriented experiments, and include a diagnostic soft-mixing ablation showing that the derived leading-tail mixing quantities can be regularized to reduce spectral leakage in practice.
\end{itemize}
